\documentclass[a4paper,12pt]{article}
\usepackage[T2A]{fontenc}
\usepackage[utf8]{inputenc}
\usepackage[russian]{babel}
\usepackage{amsmath,amssymb,mathtools}
\usepackage{iwona}
\renewcommand{\familydefault}{\rmdefault}
\usepackage{icomma}
\usepackage[a4paper,margin=19mm,top=18mm,bottom=19mm]{geometry}
\usepackage{microtype}
\usepackage{tikz}
\usetikzlibrary{calc,arrows.meta,positioning,shapes.geometric}
\usepackage{tabularx,booktabs}
\usepackage{enumitem}
\usepackage[hidelinks]{hyperref}
\usepackage{parskip}
\usepackage{fancyhdr}
\usepackage{setspace}
\usepackage{xcolor}

\definecolor{accent}{HTML}{0F6B6D}
\definecolor{darkaccent}{HTML}{07494B}
\definecolor{textgray}{HTML}{2E3133}
\definecolor{softgray}{HTML}{6B7075}
\definecolor{softaccent}{HTML}{EFF6F5}
\color{textgray}
\onehalfspacing
\setlength{\parindent}{0pt}
\setlist[itemize]{leftmargin=*,itemsep=1pt,topsep=3pt}
\setlist[enumerate]{leftmargin=*,itemsep=6pt,topsep=4pt}
\pagestyle{fancy}
\fancyhf{}
\renewcommand{\headrulewidth}{0pt}
\fancyfoot[C]{\small\color{softgray}\thepage}
\newcommand{\sectionline}{\par\vspace{1pt}\noindent\color{accent}\rule{\linewidth}{0.6pt}\par\color{textgray}\vspace{3pt}}
\newcommand{\term}[1]{\textbf{\color{darkaccent}#1}}
\newcommand{\checkitem}{\raisebox{0.1ex}{\tikz\draw[accent,line width=.65pt] (0,0) rectangle (.18,.18);}\hspace{.45em}}
\tikzset{
  flowstart/.style={ellipse,draw=accent,line width=.9pt,fill=softaccent,minimum width=37mm,minimum height=10mm,align=center,font=\small},
  flowact/.style={rounded corners=2mm,draw=accent,line width=.8pt,fill=white,text width=120mm,minimum height=11mm,align=center,font=\small,inner sep=3mm},
  flowarr/.style={-{Latex[length=2.5mm]},line width=.8pt,draw=darkaccent}
}

\begin{document}
\begin{titlepage}
\thispagestyle{empty}
\centering
\vspace*{22mm}
{\Large\color{softgray}Профильная математика · ЕГЭ\par}
\vspace{18mm}
{\fontsize{29}{35}\selectfont\bfseries\color{darkaccent}Чек-лист\par}
\vspace{7mm}
{\fontsize{21}{28}\selectfont\bfseries Логарифмы: свойства\\[2mm]и преобразование выражений\par}
\vspace{18mm}
{\large Даниил Кичук\par}
\vspace{5mm}
{\large 08.10.26\par}
\vfill
{\normalsize\color{softgray}Лёвин Артём Александрович\par}
{\small\color{softgray}эксклюзивно для Даниила Кичука\par}
\vspace{14mm}
\begin{tikzpicture}[remember picture,overlay]
  \draw[accent!55,line width=1.2pt]
  ([xshift=18mm,yshift=15mm]current page.south west)
  .. controls ([xshift=62mm,yshift=28mm]current page.south west) and ([xshift=-59mm,yshift=7mm]current page.south east) ..
  ([xshift=-18mm,yshift=18mm]current page.south east);
\end{tikzpicture}
\end{titlepage}

\section*{1. Что нужно уметь после занятия}
\sectionline
\begin{itemize}
  \item \checkitem объяснять смысл логарифма как показателя степени;
  \item \checkitem представлять числа, дроби и корни в степенном виде;
  \item \checkitem выносить показатели степени из аргумента и основания логарифма;
  \item \checkitem объединять суммы и разности логарифмов с одним основанием;
  \item \checkitem применять основное логарифмическое тождество;
  \item \checkitem менять основание логарифма и пользоваться взаимно обратными логарифмами;
  \item \checkitem комбинировать свойства степеней и логарифмов в одной задаче.
\end{itemize}

\section*{2. Смысл логарифма и ограничения}
\sectionline
\term{Логарифм} отвечает на вопрос: \emph{в какую степень нужно возвести основание $a$, чтобы получить $b$?}
\[
\boxed{\log_a b=x\ \Longleftrightarrow\ a^x=b},
\qquad a>0,\ a\ne1,\ b>0.
\]
Например, $\log_2 8=3$, поскольку $2^3=8$. Основание может быть меньше единицы: $\log_{1/2}8=-3$, поскольку $\left(\frac12\right)^{-3}=8$.

\textbf{Перед преобразованиями проверяйте:} основание положительно и отлично от $1$; аргумент строго положителен. У дробных, корневых и переменных выражений ограничения устанавливаются до упрощения.

\section*{3. Степени внутри логарифма}
\sectionline
\term{Степень в аргументе:} при $a>0$, $a\ne1$, $b>0$ и любом действительном $r$:
\[
\boxed{\log_a\!\left(b^r\right)=r\log_a b}.
\]
\term{Особый случай с модулем:} при $b\ne0$ и натуральном $n$:
\[
\boxed{\log_a\!\left(b^{2n}\right)=2n\log_a|b|}.
\]
Например, $\log_3\!\left((-3)^2\right)=\log_3 9=2\log_3|-3|=2$. Без модуля выражение $\log_3(-3)$ в действительных числах не определено.

\term{Степень в основании:} при $a>0$, $a\ne1$, $b>0$, $r\ne0$:
\[
\boxed{\log_{a^r}b=\frac{1}{r}\log_a b}.
\]
Значит, показатель из \emph{аргумента} выходит множителем $r$, а показатель из \emph{основания} --- множителем $\frac1r$.

\textbf{Пример 1.} Представим $81=3^4$, а кубический корень --- как степень $\frac13$:
\[
\log_3\!\left(\frac1{\sqrt[3]{81}}\right)
=\log_3\!\left(3^{-4/3}\right)
=-\frac43.
\]
\textbf{Пример 2.} Теперь такое же число находится в основании:
\[
\log_{\frac1{\sqrt[3]{81}}}3
=\log_{3^{-4/3}}3
=-\frac34\log_3 3=-\frac34.
\]
\textbf{Обратите внимание:} в этих примерах ответы различны, поскольку показатель расположен в разных частях логарифма.

\section*{4. Сумма, разность и коэффициент}
\sectionline
Для $a>0$, $a\ne1$, $b>0$, $c>0$:
\[
\boxed{\log_a b+\log_a c=\log_a(bc)},\qquad
\boxed{\log_a b-\log_a c=\log_a\!\left(\frac bc\right)}.
\]
Если перед логарифмом стоит множитель $k$, сначала переносим его в показатель аргумента:
\[
\boxed{k\log_a b=\log_a\!\left(b^k\right)}.
\]
\textbf{Пример.} Свернём выражение в один логарифм:
\[
\begin{aligned}
\log_2 3+2\log_2 5-\log_2 7
&=\log_2 3+\log_2 25-\log_2 7\\
&=\log_2\!\left(\frac{3\cdot25}{7}\right)
=\boxed{\log_2\frac{75}{7}}.
\end{aligned}
\]
\textbf{Ключевой порядок:} привести логарифмы к общему основанию; перенести коэффициенты в аргументы; только затем объединять сумму или разность.

\clearpage
\section*{5. Основное логарифмическое тождество}
\sectionline
Если $a>0$, $a\ne1$, $b>0$, то
\[
\boxed{a^{\log_a b}=b}.
\]
Основание внешней степени и основание логарифма должны совпадать. Полезное следствие при любом действительном $r$:
\[
\boxed{a^{r\log_a b}=b^r}.
\]
\textbf{Пример 1.} Найдём $25^{\log_{125}7}$. Представим оба основания степенями числа $5$:
\[
\begin{aligned}
25^{\log_{125}7}
&=\left(5^2\right)^{\frac13\log_5 7}
=5^{\frac23\log_5 7}
=\boxed{7^{2/3}=\sqrt[3]{49}}.
\end{aligned}
\]
\textbf{Пример 2.} Используем также свойство $a^{m+n}=a^m a^n$:
\[
\begin{aligned}
3^{1+\log_9 5}
&=3\cdot3^{\frac12\log_3 5}
=3\cdot5^{1/2}
=\boxed{3\sqrt5}.
\end{aligned}
\]
\textbf{Два полезных представления:}
\[
3=2^{\log_2 3},\qquad
5=5\log_2 2=\log_2 32.
\]
Они помогают переводить числа в удобный для задачи степенной или логарифмический вид.

\section*{6. Как менять основание логарифма}
\sectionline
Для $a>0$, $a\ne1$, $b>0$, $c>0$, $c\ne1$ справедливо:
\[
\boxed{\log_a b=\frac{\log_c b}{\log_c a}}.
\]
Частный случай --- обмен основания и аргумента:
\[
\boxed{\log_a b=\frac{1}{\log_b a}},
\qquad b\ne1.
\]
\textbf{Пример.} Логарифмы могут сократиться после перехода к одному основанию:
\[
\frac{\log_7 13}{\log_{49}13}
=\frac{\log_7 13}{\frac12\log_7 13}
=\boxed{2}.
\]
\textbf{Важно:} в формуле смены основания новое основание $c$ можно выбрать произвольно \emph{с учётом условий} $c>0$, $c\ne1$. Выбирайте число, которое облегчает преобразование.

\clearpage
\section*{Алгоритм: упрощение логарифмического выражения}
\vspace{12mm}
\begin{center}
\begin{tikzpicture}[node distance=10mm]
\node[flowstart] (s) {Начало};
\node[flowact,below=of s] (d) {Установите ОДЗ: основания положительны и не равны $1$, аргументы положительны};
\node[flowact,below=of d] (p) {Представьте подходящие числа, дроби и корни степенями одного основания};
\node[flowact,below=of p] (q) {При необходимости вынесите степени из аргументов и оснований логарифмов};
\node[flowact,below=of q] (b) {Если нужно сложить или вычесть логарифмы, сначала согласуйте их основания и коэффициенты};
\node[flowact,below=of b] (c) {Сверните суммы и разности; используйте свойства степеней и основное тождество};
\node[flowact,below=of c] (v) {Проверьте исходные ограничения и подставьте результат, если это возможно};
\node[flowstart,below=of v] (e) {Ответ};
\foreach \x/\y in {s/d,d/p,p/q,q/b,b/c,c/v,v/e}{\draw[flowarr] (\x)--(\y);}
\end{tikzpicture}
\end{center}
\clearpage

\section*{7. Несколько свойств в одном выражении}
\sectionline
\textbf{Пример 1.} Свойства степеней применяются наравне со свойствами логарифмов:
\[
\begin{aligned}
\frac{9^{\log_5 15}}{9^{\log_5 3}}
&=9^{\log_5 15-\log_5 3}
=9^{\log_5 5}
=\boxed9.
\end{aligned}
\]
\textbf{Пример 2.} В числе $1$ можно узнать логарифм с подходящим основанием:
\[
\begin{aligned}
&(1-\log_2 12)(1-\log_6 12)\\
&=(\log_2 2-\log_2 12)(\log_6 6-\log_6 12)\\
&=\log_2\frac16\cdot\log_6\frac12
=(-\log_2 6)(-\log_6 2)
=\boxed1.
\end{aligned}
\]
В последнем переходе использовано $\log_2 6\cdot\log_6 2=1$.

\section*{8. Типичные ошибки и самопроверка}
\sectionline
\begin{tabularx}{\linewidth}{@{}p{.39\linewidth}X@{}}
\toprule
\textbf{Ошибка} & \textbf{Корректный приём}\\
\midrule
$\log_a(b+c)=\log_a b+\log_a c$ & Формула суммы действует для \emph{логарифмов}, а не для сложения внутри аргумента.\\
$\log_a b\cdot\log_a c=\log_a(bc)$ & Произведение логарифмов по этой формуле не сворачивается.\\
$\log_{a^r}b=r\log_a b$ & При $r\ne0$ множитель равен $\frac1r$.\\
$\log_a((-b)^2)=2\log_a(-b)$ & Для действительных $b\ne0$ в чётной степени используйте модуль: $2\log_a|b|$.\\
$2^{\log_3 5}=5$ & Основания разные. Сначала приведите степень и логарифм к одному основанию.\\
Сложили логарифмы с разными основаниями & Сначала воспользуйтесь переходом к общему основанию.\\
\bottomrule
\end{tabularx}

\vspace{7mm}
\textbf{Контроль перед ответом}
\begin{itemize}
 \item \checkitem записаны условия существования всех исходных логарифмов;
 \item \checkitem не перепутаны показатели в аргументе и основании;
 \item \checkitem множители перенесены в степени именно своих аргументов;
 \item \checkitem сумма и разность объединены только после согласования оснований;
 \item \checkitem в основном тождестве основания действительно совпадают;
\end{itemize}

\clearpage
\section*{Тренировочный блок — Путь А}
\sectionline
\textbf{10 задач.} Выполняйте задания последовательно, записывая основные преобразования. В заданиях 5 и 6 представьте выражение \emph{одним логарифмом}; в остальных найдите числовой ответ.
\begin{enumerate}[label=\textbf{\arabic*.}]
\item Вычислите $\log_2 64$.
\item Вычислите $\displaystyle\log_3\!\left(\frac1{\sqrt[3]{81}}\right)$.
\item Вычислите $\displaystyle\log_{\sqrt[5]{2}}8$.
\item Найдите $x$, если $\displaystyle\left(\frac1{\sqrt[10]{27}}\right)^x=\sqrt[5]{81}$.
\item Сверните в один логарифм: $2\log_2 3+\log_2 5$.
\item Сверните в один логарифм: $\displaystyle\log_{\sqrt2}3+\log_4 5$.
\item Вычислите $\displaystyle25^{\log_{125}7}$. Ответ можно записать через корень.
\item Вычислите $\displaystyle3^{1+\log_9 5}$.
\item Вычислите $\displaystyle\frac{\log_7 13}{\log_{49}13}$.
\item Вычислите $\displaystyle(1-\log_2 12)(1-\log_6 12)$.
\end{enumerate}
\vfill
\textit{Самопроверка:} проверьте, что использованные свойства применимы при данных основаниях и аргументах. Не округляйте ответы, выраженные точными корнями.

\clearpage
\section*{Ответы}
\sectionline
\renewcommand{\arraystretch}{1.42}
\begin{tabularx}{\linewidth}{@{}>{\bfseries}c X@{}}
\toprule
№ & Ответ\\
\midrule
1 & $6$\\
2 & $-\dfrac43$\\
3 & $15$\\
4 & $x=-\dfrac83$\\
5 & $\log_2 45$\\
6 & $\log_2(9\sqrt5)$\\
7 & $\sqrt[3]{49}$\\
8 & $3\sqrt5$\\
9 & $2$\\
10 & $1$\\
\bottomrule
\end{tabularx}
\vspace{11mm}
\textbf{Ключевые шаги для самопроверки}
\begin{itemize}
 \item Задания 2--4: $81=3^4$, $27=3^3$, $\sqrt[5]2=2^{1/5}$.
 \item Задания 5--6: сначала переносите коэффициенты, затем применяйте формулу суммы.
 \item Задания 7--8: приводите степени и логарифмы к единому основанию.
 \item Задание 9: $49=7^2$, поэтому $\log_{49}13=\frac12\log_7 13$.
 \item Задание 10: замените $1$ на $\log_2 2$ и $\log_6 6$, затем используйте взаимно обратные логарифмы.
\end{itemize}
\vfill
\begin{center}
\small\color{softgray}Лёвин Артём Александрович · эксклюзивно для Даниила Кичука · 08.10.26
\end{center}
\end{document}